\documentclass[11pt,a4paper]{article}
\usepackage[margin=1in]{geometry}
\usepackage[T1]{fontenc}
\usepackage[utf8]{inputenc}
\usepackage{lmodern}
\usepackage{microtype}
\usepackage{hyperref}
\usepackage{booktabs}
\usepackage{tabularx}
\usepackage{array}
\usepackage{amsmath,amssymb}
\usepackage{xcolor}
\usepackage{url}

\hypersetup{
  colorlinks=true,
  linkcolor=blue!50!black,
  citecolor=blue!50!black,
  urlcolor=blue!50!black,
  pdftitle={How Realistic Are Successful Evasion Attacks?},
  pdfauthor={Sathyaraj Kolandasamy}
}

\title{How Realistic Are ``Successful'' Evasion Attacks?\\
\large A Constraint-Aware Reproducible Comparison of Adversarial Attacks\\
Against ML-Based Network Intrusion Detection}
\author{Sathyaraj Kolandasamy\\
Nova Southeastern University\\
College of Computing, AI, and Cybersecurity\\
\texttt{https://github.com/sathya-ksamy-nsu/constraint-aware-nids}}
\date{Preprint v0.1 \quad 25 August 2026}

\begin{document}
\maketitle

\begin{abstract}
Machine-learning-based network intrusion detection systems (NIDS) are widely evaluated under adversarial evasion, and many papers report very high attack success rates (ASR). Surveys from 2023--2024 warn that those successes are often \textbf{unrealistic}: attacks perturb flow features independently, violating protocol semantics and feature interdependencies that a packet-sending adversary cannot break. This preprint contributes a \textbf{constraint-aware evaluation protocol} and an \textbf{open harness} to measure that gap. We formalize a reusable \emph{realistic-perturbation mask} (mutability, direction, range, and interdependency constraints), compare gradient-based (FGSM/PGD), decision-based (HopSkipJump), and generative (GAN) attacks in unconstrained vs.\ constrained regimes, and report ASR together with $L_2$/$L_\infty$ size and a \textbf{valid-sample rate}. The hypothesis, to be tested on CICIDS-2017 or UNSW-NB15, is that enforcing semantically valid, functionality-preserving perturbations \textbf{lowers reported ASR and can reorder attack/defense rankings}. A synthetic end-to-end run of the harness shows the intended mechanism: unconstrained HopSkipJump reached ASR~$=1.00$ with valid-sample rate~$=0.00$, while the constrained regime forced valid-rate~$=1.00$ and collapsed ASR to $0.027$ (undefended MLP). Those figures are \textbf{pipeline validation, not NIDS evidence}. The artifact is the protocol, mask, and one-command matrix---intended as a citable baseline for follow-on work and for a forthcoming workshop version (AI-SEC 2027).
\end{abstract}

\noindent\textbf{Intended arXiv categories:} \texttt{cs.CR} (primary), \texttt{cs.LG} (cross-list).\\
\textbf{Code (MIT):} \url{https://github.com/sathya-ksamy-nsu/constraint-aware-nids}

\paragraph{Preprint status.}
This is a Phase~1 arXiv preprint of an evaluation \emph{protocol and open harness}. Benchmark results on CICIDS-2017 / UNSW-NB15 are \textbf{not yet reported}. Section~\ref{sec:results} includes a pipeline-validation run on \textbf{synthetic data only}; those numbers are \textbf{not} findings about real intrusion detectors.

\section{Introduction}

Modern networks overwhelm signature-based intrusion detection. ML-based NIDS learn statistical regularities of benign and malicious flows and are now common in research and production stacks \cite{he2023aml}. Their decision boundary is therefore a security surface, not only an accuracy metric.

Adversarial examples show that small perturbations can flip a classifier \cite{goodfellow2015explaining,madry2018towards}. Against a NIDS, evasion means malicious traffic is labeled benign. The literature reports many such attacks with high ASR \cite{he2023aml,sharma2024systematic}.

A structural flaw remains: most attacks operate in \textbf{feature space}, treating each CICFlowMeter-style coordinate as independently mutable under a norm budget. A flow vector is a \emph{derived summary} of packets. An attacker who can only send packets cannot independently set inter-arrival statistics, durations, and byte totals to inconsistent values \cite{alhussien2024constraining,ennaji2024adversarial}. Unconstrained ASR can therefore overstate the threat, and defenses can be tuned against impossible traffic.

\paragraph{Research question.}
When perturbations against an ML-based NIDS are restricted to semantically valid, functionality-preserving changes, how do ASR and the ranking of attacks and defenses change relative to unconstrained feature-space evaluation?

\paragraph{Hypothesis.}
Constraining perturbations \textbf{lowers reported ASR} and \textbf{reorders} attack/defense rankings versus unconstrained perturbations.

\paragraph{Contributions.}
\begin{enumerate}
\item A formal, reusable \textbf{constraint specification} (realistic-perturbation mask).
\item A \textbf{reproducible harness} (identical preprocessing/splits, ART-based attacks, valid-sample accounting).
\item A controlled unconstrained-vs-constrained \textbf{experiment design} spanning white-box, black-box, and generative attacks, with and without adversarial training.
\item A related-work table annotated by whether prior attacks respect domain constraints.
\item An \textbf{open implementation} with unit tests and a documented synthetic pipeline-validation run.
\end{enumerate}

\section{Background and threat model}

\subsection{ML-based NIDS pipeline}
Packets are grouped into flows (five-tuple). Feature extraction (e.g.\ CICFlowMeter) yields counts, durations, flag tallies, and rates. After preprocessing, a classifier outputs benign/malicious. Not every point in feature space is the image of a realizable packet stream.

\subsection{Attack families}
\textbf{FGSM} \cite{goodfellow2015explaining}: single-step sign of the loss gradient; baseline.
\textbf{PGD} \cite{madry2018towards}: iterated projected steps; strong first-order white-box attack.
\textbf{HopSkipJump / Boundary / ZOO:} query-based black-box attacks; applicable to non-differentiable models such as random forests.
\textbf{GAN-based evasion:} generate evasive vectors whose distribution resembles benign traffic.

\subsection{Knowledge}
White-box (gradients), grey-box (surrogate/transfer), black-box (queries only). A deployed NIDS is often closest to black-box; white-box remains the worst-case stress test.

\subsection{Feature space vs.\ problem space}
Feature-space attacks edit the vector. Problem-space attacks edit packets; features change only as a consequence. This protocol is \textbf{constrained feature-space}: a projection/filter approximating problem-space legality without yet crafting packets. Valid-sample rate makes residual unrealizability visible. True packet-level validation is future work.

\subsection{What ``realistic'' means here}
A perturbation must be \textbf{functionality-preserving}, \textbf{protocol-valid}, and \textbf{interdependency-consistent}. Typical rules: immutable protocol/port/flag-type identifiers; non-negative, often increase-only counts; $\mathrm{total\_packets}=\mathrm{fwd}+\mathrm{bwd}$; rates consistent with counts and duration.

\section{Related work}

He et al.\ \cite{he2023aml} survey adversarial ML for NIDS and flag the realizability mismatch. Ennaji et al.\ \cite{ennaji2024adversarial} catalogue domain constraints and the difficulty of mapping feature-space $\delta$ back to valid traffic. Sharma and Chen \cite{sharma2024systematic} document inconsistent evaluation protocols. Alhussien et al.\ \cite{alhussien2024constraining} show that constraining attacks changes transferability and defense conclusions. This preprint complements that line by releasing a mask and harness and a paired unconstrained/constrained matrix on identical splits.

\begin{table}[t]
\centering
\caption{Prior attacks vs.\ domain constraints.}
\label{tab:related}
\small
\begin{tabular}{p{3.4cm}p{2.4cm}p{1.8cm}p{2.6cm}p{2.4cm}}
\toprule
Attack & Threat model & Space & Realistic constraints? & Notes \\
\midrule
FGSM \cite{goodfellow2015explaining} & White-box & Feature & No (default) & Single-step baseline \\
PGD \cite{madry2018towards} & White-box & Feature & No (default); mask can be applied & Robustness standard \\
HopSkipJump / Boundary & Black-box (decision) & Feature & No (default) & Works on RF \\
ZOO & Black-box (scores) & Feature & No (default) & Query-expensive \\
GAN-based evasion & Grey/black-box & Often feature & Partial & Distribution-aware \\
Constrained attacks \cite{alhussien2024constraining} & Mixed & Feature + domain & Yes & Closest prior \\
\bottomrule
\end{tabular}
\end{table}

\paragraph{Gap.}
Unconstrained feature-space ASR dominates published tables. Missing is a shared, released protocol that reports \textbf{valid-sample rate} next to ASR so follow-on papers can cite a common baseline.

\section{Evaluation protocol}

\subsection{Data (intended benchmark)}
Use \textbf{one} public dataset with identical preprocessing for every cell: CICIDS-2017 \cite{sharafaldin2018cicids} (preferred for constraint provenance) or UNSW-NB15 \cite{moustafa2015unsw}. Raw files are \textbf{not} redistributed. Splits: stratified train/val/test with recorded seed (\texttt{config.yaml}: seed 42; val/test $0.15$ each). Standardization is fit on train only.

\subsection{Models}
\textbf{MLP} (PyTorch, ART-wrapped) for white-box FGSM/PGD.
\textbf{Random forest} (scikit-learn) for black-box HopSkipJump/GAN only.

\subsection{Attacks and defense}
Implemented with the Adversarial Robustness Toolbox (FGSM, PGD, HopSkipJump, optional ZOO) plus a simple GAN attack. Each attack is run unconstrained and constrained (mask projection + validity filter). Defense: PGD adversarial training \cite{madry2018towards}.

\subsection{Constraint mask (central artifact)}
For $x'=x+\delta$, require:
\begin{itemize}
\item \textbf{Mutability:} immutable features have $\delta_i=0$ (protocol, many ports/flags).
\item \textbf{Direction:} e.g.\ increase-only packet/byte counts.
\item \textbf{Range:} non-negativity and protocol boxes.
\item \textbf{Interdependencies:} e.g.\ total packets $=$ forward $+$ backward; ratios repaired after projection.
\end{itemize}
The implementation (\texttt{src/constraints.py}) projects then scores valid-sample rate.

\subsection{Metrics and matrix}
\textbf{ASR}: fraction of originally detected malicious samples that become benign after perturbation.
Mean $L_2$ / $L_\infty$ of $\delta$.
\textbf{Valid-sample rate}: fraction of adversarial samples satisfying the mask.

Matrix: $\{\mathrm{FGSM},\mathrm{PGD},\mathrm{HopSkipJump},\mathrm{GAN}\} \times \{\mathrm{unconstrained},\mathrm{constrained}\} \times \{\mathrm{no\ defense},\mathrm{adv.\ training}\} \times \{\mathrm{MLP},\mathrm{RF}\ \mathrm{(black\text{-}box\ only)}\}$.

\section{Experimental setup}

Pipeline-validation run (31 July 2026): Python 3.13.3, PyTorch 2.10.0+cpu, ART 1.20.1, scikit-learn 1.8.0, Windows 10, CPU. Unit tests: 27 passed.

\begin{verbatim}
pip install -r requirements.txt
python -m pytest tests/ -v
python experiments/run_experiment.py --synthetic --model mlp --attack all
python experiments/run_experiment.py --model mlp --attack all
\end{verbatim}

\section{Results}
\label{sec:results}

\subsection{Pipeline validation (synthetic data --- not a NIDS finding)}

Gaussian-blob synthetic features, MLP, full 16-cell matrix (\texttt{RESULTS.md}). Clean test accuracy $\approx 0.993$.

\begin{table}[t]
\centering
\caption{Undefended MLP on synthetic data (pipeline validation only).}
\label{tab:synth}
\small
\begin{tabular}{llrrr}
\toprule
Attack & Regime & ASR & Valid-rate & Mean $L_2$ \\
\midrule
FGSM & unconstrained & 0.000 & 0.327 & 0.346 \\
FGSM & constrained & 0.000 & 1.000 & 0.480 \\
PGD & unconstrained & 0.000 & 0.327 & 0.346 \\
PGD & constrained & 0.000 & 1.000 & 0.480 \\
HopSkipJump & unconstrained & \textbf{1.000} & \textbf{0.000} & 2.923 \\
HopSkipJump & constrained & \textbf{0.027} & \textbf{1.000} & 2.325 \\
GAN & unconstrained & 0.000 & 0.380 & 0.133 \\
GAN & constrained & 0.000 & 1.000 & 0.325 \\
\bottomrule
\end{tabular}
\end{table}

Unconstrained HopSkipJump ``succeeds'' on every attacked malicious sample \textbf{and none of those samples are valid under the mask}. Constraining forces valid-rate $=1$ and ASR falls to $0.027$. Gradient attacks did not evade this easy synthetic MLP at the configured $\varepsilon$ (ASR $=0$); that is a limitation of the smoke-test distribution, not evidence that FGSM/PGD are weak on CICIDS. RF black-box cells were not completed in that run.

\subsection{Intended CICIDS-2017 / UNSW-NB15 matrix (pending)}
All CICIDS/UNSW cells are TBD pending a real-dataset run. Predicted pattern (to confirm or refute): constrained ASR $\ll$ unconstrained ASR; unconstrained valid-rate low; ranking may reorder.

\section{Discussion}
If the CICIDS/UNSW gap matches the synthetic \emph{mechanism} (high unconstrained ASR, near-zero valid-rate), then many published NIDS ASR figures are not operational threats. Defenders should demand valid-sample rate next to ASR and should measure adversarial training under the same mask the network actually imposes.

\section{Limitations}
Single-dataset scope once chosen; mask may be incomplete or overly strict; constrained feature-space $\neq$ packet crafting; synthetic Section~\ref{sec:results} does not transfer; attack/defense coverage is three families plus one defense; RF matrix incomplete.

\section{Ethics}
Evasion research is framed \textbf{defensively}: improve evaluation of detectors. No exploit payloads. Datasets remain with their providers. Do not treat synthetic ASR as evidence against production NIDS.

\section{Conclusion and next versions}
We release a constraint-aware adversarial-NIDS protocol and harness, and we document a synthetic run that exercises valid-sample accounting. Version 0.2 should replace the pending table with CICIDS-2017 or UNSW-NB15 measurements, then serve as the arXiv record cited by the AI-SEC 2027 workshop submission (target deadline 10 December 2026).

\section*{Reproducibility}
Code: \url{https://github.com/sathya-ksamy-nsu/constraint-aware-nids}.
Config/seed: \texttt{config.yaml} (\texttt{seed: 42}).
Mask: \texttt{src/constraints.py}.
Pipeline-validation log: \texttt{RESULTS.md}.
Data download: \texttt{data/README.md} (do not redistribute raw CIC/UNSW files).

\bibliographystyle{plain}
\begin{thebibliography}{99}

\bibitem{he2023aml}
K.~He, D.~S.~Kim, and M.~R.~Asghar.
Adversarial machine learning for network intrusion detection systems: A comprehensive survey.
\emph{IEEE Communications Surveys \& Tutorials}, 25(1):538--566, 2023.

\bibitem{alhussien2024constraining}
N.~Alhussien, A.~Aleroud, A.~Melhem, and S.~Y.~Khamaiseh.
Constraining adversarial attacks on network intrusion detection systems: Transferability and defense analysis.
\emph{IEEE Transactions on Network and Service Management}, 21(3):2751--2772, 2024.

\bibitem{sharma2024systematic}
S.~Sharma and Z.~Chen.
A systematic study of adversarial attacks against network intrusion detection systems.
\emph{Electronics}, 13(24):5030, 2024.

\bibitem{ennaji2024adversarial}
S.~Ennaji, F.~De~Gaspari, D.~Hitaj, A.~K.~Bidi, and L.~V.~Mancini.
Adversarial challenges in network intrusion detection systems: Research insights and future prospects.
arXiv:2409.18736, 2024.

\bibitem{madry2018towards}
A.~Madry, A.~Makelov, L.~Schmidt, D.~Tsipras, and A.~Vladu.
Towards deep learning models resistant to adversarial attacks.
In \emph{ICLR}, 2018.

\bibitem{goodfellow2015explaining}
I.~J.~Goodfellow, J.~Shlens, and C.~Szegedy.
Explaining and harnessing adversarial examples.
In \emph{ICLR}, 2015.

\bibitem{sharafaldin2018cicids}
I.~Sharafaldin, A.~H.~Lashkari, and A.~A.~Ghorbani.
Toward generating a new intrusion detection dataset and intrusion traffic characterization.
In \emph{ICISSP}, 2018.

\bibitem{moustafa2015unsw}
N.~Moustafa and J.~Slay.
UNSW-NB15: A comprehensive data set for network intrusion detection systems.
In \emph{MilCIS}, 2015.

\end{thebibliography}

\end{document}
